\section{Mean Curvature Flow}

The mean curvature flow is (after choosing an appropriate `gauge') a parabolic equation, with the minimal surface equation as its static case and elliptic counterpart. Its linear analogue is the heat equation, with its elliptic counterpart being the Laplace equation.

\begin{definition}[Mean Curvature Flow]\label{def:Mean Curvature Flow}
    A smooth 1-parameter family of immersions $X:M^n\times I\to\mathbb{R}^{n+1}$ satisfies the mean curvature flow if
    \begin{equation}\label{eq:M-0-1}
        \partial_t X(\boldsymbol{x},t)=\boldsymbol{H}(\boldsymbol{x},t),\,\,\,\,\,\,\,\forall (\boldsymbol{x},t)\in M^n\times\left[0,T\right)
    \end{equation}
    where $\boldsymbol{H}$ denotes the mean curvature vector of the family.
\end{definition}

We say that $X:M\times I\to\mathbb{R}^{n+1}$ is a solution of \textbf{reparametrized mean curvature
flow} if
\begin{equation*}
    {\left[\partial_t X(\boldsymbol{x},t)\right]}^{\bot}=\boldsymbol{H}(\boldsymbol{x},t)
\end{equation*}
In this case, the vector field
\begin{equation*}
    V(\boldsymbol{x},t):={\left[\partial_t X(\boldsymbol{x},t)\right]}^{\bot}-\partial_t X(\boldsymbol{x},t)
\end{equation*}
is tangential. Defining the 1-parameter family of self-diffeomorphisms $\varphi_t:M\to M$ by
\begin{equation}\label{eq:M-0-2}
    \begin{cases}
        \partial_t\varphi_t(\boldsymbol{x})=dX^{-1}_t[V(\varphi_t(\boldsymbol{x}),t)],\\
        \varphi_0=\mathrm{id}_M
    \end{cases}
\end{equation}
and let $\hat{X}_t=X_t\circ\varphi_t$. Then we claim that $\hat{X}$ is a solution of the mean curvature flow (\ref{eq:M-0-1}). This is because
\begin{align*}
    \partial_t\hat{X}(\boldsymbol{x},t)&=\left[\frac{\partial}{\partial\boldsymbol{x}}X_t,\frac{\partial}{\partial t}X_t\right]\cdot\left[\begin{array}{c}
        \partial_t\varphi_t(\boldsymbol{x})\\
        \partial_t t
    \end{array}\right]=dX_t(\varphi_t)\cdot\partial_t\varphi_t(\boldsymbol{x})+\partial_t X_t(\varphi_t)
\end{align*}
Combined with (\ref{eq:M-0-2}) and the relationship between differentiation and the action of tangent vectors, we can obtain that
\begin{align*}
    dX_t(\varphi_t)\cdot\partial_t\varphi_t(\boldsymbol{x})&=dX_t[dX^{-1}_t[V(\varphi_t,t)]]\\
    &=V(\varphi_t,t)={\left[\partial_t X(\varphi_t,t)\right]}^{\bot}-\partial_t X(\varphi_t,t)
\end{align*}
Hence, by the condition of the reparametrized mean curvature flow we have that
\begin{equation*}
    \partial_t \hat{X}(\boldsymbol{x},t)={\left[\partial_t X(\varphi_t,t)\right]}^{\bot}=\boldsymbol{H}(\varphi_t,t)=\boldsymbol{H}_{\hat{X}}(\boldsymbol{x},t)
\end{equation*}

A much easier solution of the MCF (Mean Curvature Flow) is the following shrinking round sphere. Denote by $S^n_r\subseteq\mathbb{R}^{n+1}$ the $n$-sphere of radius $r>0$ centered at the origin. Set $\mathcal{M}_0=S^n_{r_0}$ where $r_0>0$. We seek a solution of a family of $n$-sphere to mean curvature flow, starting at $\mathcal{M}_0$. Then let $\mathcal{M}_t=S^n_{r(t)}$, it can naturally parameterized by an embedding $X_t:S^n_1\hookrightarrow $ as
\begin{equation*}
    X_t:\boldsymbol{x}\mapsto r(t)\cdot\boldsymbol{x}
\end{equation*}
Then by Example.\ref{exam:H-1-1} we could calculate its mean curvature and $\partial_t X(\boldsymbol{x},t)$ as
\begin{equation*}
    \boldsymbol{H}=-H(\boldsymbol{x},t)\cdot\mathbf{N}(\boldsymbol{x},t)=-n\frac{\boldsymbol{x}}{r(t)},\,\,\,\, \partial_t X(\boldsymbol{x},t)=\frac{\mathrm{d}r(t)}{\mathrm{d}t}\boldsymbol{x}
\end{equation*}
Then the MCF (\ref{eq:M-0-1}) relates to an ODE
\begin{equation*}
    \frac{\mathrm{d}r(t)}{\mathrm{d}t}=-\frac{n}{r(t)}
\end{equation*}
One can solve this equation with initial value $r(0)=r_0$ as
\begin{equation*}
    r(t)=\sqrt{r^2_0-2nt}
\end{equation*}
Then we have
\begin{equation*}
    X_t(\boldsymbol{x})=\boldsymbol{x}\cdot\sqrt{r^2_0-2nt}
\end{equation*}
Which means that the sphere in such family would shrink to a point for $t\in (-\infty,r^2_0/2n)$

Recall that a minimal hypersurface $\mathcal{M}^n\subseteq\mathbb{R}^{n+1}$ is a hypersurface satisfying $H\equiv 0$. Given a minimal surface $\mathcal{M}^n_0$ its clear that $\mathcal{M}_t\equiv\mathcal{M}_0$ is a solution of the mean curvature flow. In the following section we will discuss it specifically.

We claim that if ${\mathcal{M}^{n-m}_t}$ is a solution to mean curvature flow in $\mathbb{R}^{n-m+1}$, then the family of product hypersurfaces
\begin{equation*}
    \mathcal{M}^n_t :=\mathcal{M}^{n-m}_t\times\mathbb{R}^m\subseteq\mathbb{R}^{n+1}
\end{equation*}
should be a solution to mean curvature flow in $\mathbb{R}^{n+1}$ which is called the cylinder over a solution. To see this let the immersion over $\mathcal{M}^{n-m}_t$ be $X^\prime(\boldsymbol{x},t)$ which satisfies that
\begin{equation*}
    \partial_t X^\prime(\boldsymbol{x},t)=\boldsymbol{H}_{n-m}
\end{equation*}
Then let the immersion $X:\mathcal{M}^n_t\to\mathbb{R}^{n+1}$ be defined as $X(\boldsymbol{x},\boldsymbol{y},t)=(X^\prime(\boldsymbol{x},t),\boldsymbol{y})$. Let $\partial/\partial x_i, i=1,2,\dots,n-m$ denote the linearly independent tangent vectors of $\mathcal{M}^{n-m}_t$ and $\partial/\partial y_j, j=1,2,\dots,m$ denote the tangent vector in $\mathbb{R}^m$. Thus we have
\begin{equation*}
    g^X_{ij}=\left[\begin{matrix}
        g^{X^\prime}_{i_1 j_1}(\boldsymbol{x},t)&    0\\
        0&    \delta_{i_2 j_2}\\
    \end{matrix}\right]
\end{equation*}
where $\delta_{i_2 j_2}=1$ if and only if $i_2=j_2$ and $g^{X^\prime}_{i_1 j_1}(\boldsymbol{x},t)$ is the first fundamental form of $\mathcal{M}^{n-m}_t$ over $X^\prime$. With this its easy to verify that the second fundamental form of $\mathcal{M}^n_t$ should be
\begin{equation*}
    \mathbf{II}=h^X_{ij}=\left[
        \begin{matrix}
            h^{X^\prime}_{i_1,j_1}(\boldsymbol{x},t)&    0\\
            0&   0\\
        \end{matrix}
    \right]
\end{equation*}
Hence we can get that
\begin{equation*}
    \boldsymbol{H}_{n}=-g^{ij}_X\cdot h^X_{ij}=\left(\boldsymbol{H}_{n-m},0\right)
\end{equation*}
Therefore, we deduce that
\begin{equation*}
    \frac{\partial}{\partial t}X(\boldsymbol{x},\boldsymbol{y},t)=\left(\partial_t X^\prime (\boldsymbol{x},t),0 \right)=(\boldsymbol{H}_{n-m},0)=\boldsymbol{H}_n
\end{equation*}

\subsection{Minimal Surface Problem}
Now, let us briefly review the origins of the concept of mean curvature flow. First and formost, let's consider a natural problem as follows, which is called \textbf{Minimal Surface Problem}. In $\mathbb{R}^3$, it can be expressed as:
\begin{problem}[Plateau]\label{prob:Plateau}
    Given a closed curve $\Gamma\in\mathbb{R}^3$, find a surface $M$ with boundary $\Gamma$ such that its area is minimal among all surfaces bounded by $\Gamma$. 
\end{problem}

To find such surface $M$ let it be parameterized by $X:\Omega\to\mathbb{R}^3$ where the parameters $(u,v)\subseteq\Omega$ and $X\mid_{\partial\Omega}=\Gamma$ then its area can be written as
\begin{equation*}
    A_M=\int_\Omega\left|X_u\times X_v\right|\,\mathrm{d}u\,\mathrm{d}v=\int_{\Omega}\sqrt{EG-F^2}\, \mathrm{d}\sigma
\end{equation*}
where $\mathrm{d}\sigma=\mathrm{d}u\wedge\mathrm{d}v$ and $E=\left<X_{u},X_u\right>$, $F=\left<X_u,X_v\right>$, $G=\left<X_v,X_v\right>$. Then consider a deformation $X_t=X(u,v,t)=X+t\phi\cdot\mathbf{N}$ where $\phi$ is a compact support function $\phi:\Omega\to\mathbb{R}$ such that $\phi\mid_{\partial\Omega}=0$ and $\mathbb{N}$ is the unit normal vector field of $M$. Then we know the first order variation of it should be
\begin{equation*}
    \delta(A_{X_t})={\left.\frac{\mathrm{d}}{\mathrm{d}t}A_{X_t}\right|}_{t=0}=-\int_{\Omega}H\cdot\phi\sqrt{g}\,\mathrm{d}\sigma=-\int_M H\cdot\phi\,\mathrm{d}A
\end{equation*}
Under the inner product $\left<f,g\right>=\int_{M}f\cdot g\,\mathrm{d}A$ in $L^2(M)$ we know that
\begin{equation*}
    \left<-H,\phi\right>=\delta\left(A_{X_t}\right)=\left<\nabla A,\phi\right>
\end{equation*}
Which shows $\nabla A=-H$. Thus the mean curvature flow of any $X:\Omega\times I\to\mathbb{R}^3$
\begin{equation*}
    \begin{cases}
        \partial_t X(u,v,t)=\boldsymbol{H}=H\cdot\mathbf{N}=-\nabla A\\
        X\mid_{\partial\Omega}(u,v,t)=\Gamma, &\forall\, (u,v)\in\Omega, t\in I\\
    \end{cases}
\end{equation*}
represent a family of surfaces bounded by $\Gamma$ such that their area decreased. Since $M$ is the minimal surface, then by Fermat lemma its area variation should always be zero, that is
\begin{equation*}
    \delta(A_{X_t})=\left<\nabla A,\phi\right>\equiv 0,\,\,\, \forall\phi
\end{equation*}
It implies that $\nabla A=-H\cdot\mathbf{N}\equiv 0$, then $H\equiv 0$. Thus, each minimal surface should necessarily be a stable point of a mean curvature flow. Particularly, in $\mathbb{R}^3$ we have that
\begin{equation*}
    H=\mathrm{tr}(\mathrm{L})=\mathrm{tr}(g^{ij}\cdot h_{ij})=\frac{EN-2FM+GL}{EG-F^2}
\end{equation*}
where $L=\left<X_{uu},\mathbf{N}\right>$, $M=\left<X_{uv},\mathbf{N}\right>$, $N=\left<X_{vv},\mathbf{N}\right>$. Then by $H=0$ we can deduce the equation of a minimal surface $z=f(x,y)$ with the boundary $\Gamma: I\to\mathbb{R}^3$ should be
\begin{equation}\label{eq:M-1-1}
    \begin{cases}
        {\left[1+{(f^\prime_y)}^2\right]}f_{x}^{''}-2f^\prime_x f^\prime_y f^{''}_{xy}+{\left[1+{(f^\prime_x)}^2\right]}f^{''}_{yy}=0\\
        \Gamma_z(t)=f\left[\Gamma_x(t),\Gamma_y(t)\right]\\
    \end{cases}
\end{equation}

\begin{example}\label{exam:M-1-1}
    The plane and the following two surfaces are the classical solution of the Plateau's problem.
    \begin{enumerate}[(i)]
        \item For any closed simple curve $\Gamma$ on the following Catenoid, the Catenoid itself should be a minimal surface of such Plateau's problem with boundary $\Gamma$.
        \begin{equation}\label{eq:M-1-2}
            X(u,v)=\left(u,\cosh{u}\cos{v},\cosh{u}\sin{v}\right)
        \end{equation}
        \item For any closed simple curve $\Gamma$ on the following Helicoid, it is a minimal surface of such Plateau's problem with boundary $\Gamma$.
        \begin{equation}\label{eq:M-1-3}
            X(u,v)=\left(v\cos{u},v\sin{u},u\right)
        \end{equation}
    \end{enumerate}
\end{example}

One can verify this example by simply checking the mean curvature $H$ of the surface (\ref{eq:M-1-2}) and (\ref{eq:M-1-3}) are both zero.

Furthermore, based on the concept of hypersurface area established in the previous section, we can generalize Plateau's problem to arbitrary $n$-dimensional hypersurfaces and demonstrate that a necessary condition for the solutions to these problems is that they constitute a stationary point of the mean curvature flow (we shall prove this by verifying that the mean curvature flow is the gradient flow of the area functional).
\begin{problem}[Minimal Surface]\label{prob:M-1-2}
    For a smooth, compact and closed $(n-1)$-dimensional submanifold $\Gamma\subseteq \mathbb{R}^{n+1}$ find an $n$-dimensional hypersurface $M$ which could be bounded by $\Gamma$ and has the minimal area among all the hypersurface with the same boundary.
\end{problem}

\subsection{Isoperimetric Problem}

\begin{problem}[Isoperimetric Problem]\label{prob:M-2-1}
    Given an area $A>0$ in the plane $\mathbb{R}^2$, does there exist a closed curve with the shortest possible parameter such that the area enclosed within it is exactly $A$?
\end{problem}

We claim that such closed curve does exist, because the area enclosed by a closed curve of zero length is evidently also zero, it follows that the length of any closed curve enclosing an area $A>0$ must necessarily be greater than zero. Furthermore, as it can be proven that the set comprising all closed curves with an enclosed area of $A$ is complete, the Extreme Value Theorem implies that the lengths of all curves within this set must possess a greatest lower bound.

To find such curve we define the functional of the length of any closed curve $\gamma:I\to\mathbb{R}^2$ as
\begin{equation*}
    L_\gamma=\int_I\left|\dot{\gamma}\right|\,\mathrm{d}t
\end{equation*}

\subsection{Self-similar Solutions}